\SourceLeaf{285}%
\(GF\) at \(G\)). Denote likewise the (always negative) values of \(\tfrac{206265''\mathrm{d}m}{2m\mathrm{d}S}\) at those same points by \(-k^0,-k'\). Then
\[
\begin{aligned}
206265''l^0&=-2k^0\sin V^0\\
206265''l'&=+2k'\sin V'
\end{aligned}
\]

When converted into seconds, the expressions for \(\psi^0,\psi'\) given in the preceding article therefore become, if we omit the factors \(\BareRoot[6]{\tfrac{m^0}{m'}},\BareRoot[6]{\tfrac{m'}{m^0}}\) that differ only imperceptibly from unity here,
\[
\begin{aligned}
\psi^0&=-\tfrac{1}{3}h(2k^0\sin V^0-k'\sin V')\\
\psi'&=-\tfrac{1}{3}h(2k'\sin V'-k^0\sin V^0)
\end{aligned}
\]

In its last column, headed \(k\), the table appended to this memoir gives the values of \(k^0,k'\) for the corresponding values of \(S\), which are to be found in the first column under the heading \(Q+q\). Since \(k\) is always positive and \(\sin V^0,\sin V'\) always have opposite signs, \(\psi^0\) is negative and \(\psi'\) positive when \(G\) lies west of \(F\), and conversely. In making the calculation, remember that in these formulas \(h\) is understood to be expressed in parts of the radius. The distance between the points \(F,G\), given in any unit of length, must therefore first be divided by the value of \(A\) expressed in the same unit.

\begingroup\rightskip=0pt plus 1em
In our conformal mapping from the ellipsoid to the sphere, a meridian on the former is again represented by a meridian on the latter. It is therefore clear that each element of \(L\) has the same inclination to the meridian as the corresponding element of \(M\), and hence that the azimuths of the geodesic at its two endpoints will be \(V^0+\psi^0\) and \(V'+\psi'\), respectively. Conversely, if these are given, they are reduced to the sphere by applying \(-\psi^0,-\psi'\). In calculating these corrections, which are always nearly zero, it clearly makes no difference if the azimuths on the ellipsoid are used in the above formulas in place of \(V^0,V'\).\par\endgroup

\Article{14}
To calculate the reductions of the directions according to the rules given, for transfer from the ellipsoid to the sphere or vice versa,
